\scanpage{49}
$\unsure{R|\mathcal{R}}$ non-\unsure{mixed|pointwise} \unsure{equivalence|equivalence} relation on $S$.

\inserted{$\underline{\mathcal{R}}$}-\uline{saturated} \uline{segments} : \unsure{stable|stable|stable} under passage to the \unsure{complement|complementation}, and under intersection and union when these \deleted{\unsure{d}}give segments (\inserted{for example if} $A\cup B\ne S$, $A\cap B\ne\phi$)

If $S_i\in\mathcal{R}$, then the component segments of $\complement S_i$ are saturated, hence also their complements i.e. the minimal segments containing $S_i$ (or cocomponent \inserted{segments} of $S_i$). \deleted{More \unsure{generally|generally}, \unsure{If|if} $T\subset S$ is an $R$-saturated subset of $S$, then the component segments of}

\noindent\llap{\rule[-.92\baselineskip]{.4pt}{1.8\baselineskip}\hspace{.35em}}\uline{Proposition} The minimal saturated segments are the $\deleted{\unsure{A|S}}S_i\in R$ which are segments.

% slanted
\reversemarginpar\margin{One calls the subsets \uline{\unsure{tabs|little angles}}}

It is evident that if $S_i\in\deleted{\unsure{A}}R$ is a segment, then it is a saturated segment, and it is minimal among the latter (and even among the nonempty saturated subsets). Conversely, let $A$ be a minimal saturated segment, \deleted{\unsure{let}} let $S_i\in R$, $S_i\subset A$, \deleted{let} \unsure{I|I} say that $S_i=A$, \deleted{\unsure{or|sin}} i.e. that one cannot have $S_i,S_j\in R$ distinct, and
% slanted
\margin{\unsure{omission|forgotten} on the side of the \unsure{question|section|notion|notation}}
\par\noindent\begin{minipage}[t]{.23\linewidth}
\vspace{0pt}
\figure{\begin{tikzpicture}[x=1cm,y=1cm,line cap=round,line join=round]
  \draw[line width=.5pt] (.18,0) .. controls (-.10,.48) and (-.10,1.13) .. (.15,1.67)
    .. controls (.30,2.10) and (.79,2.48) .. (1.08,2.72);
  \draw[line width=2.2pt] (.15,1.67) .. controls (.30,2.10) and (.79,2.48) .. (1.08,2.72);
  \fill (.18,0) ellipse [x radius=.049,y radius=.037];
  \fill (1.08,2.72) ellipse [x radius=.044,y radius=.052];
\end{tikzpicture}}
\end{minipage}\hfill\begin{minipage}[t]{.73\linewidth}
\vspace{0pt}
$\subset A$. Indeed, if $A_i$ $(A_j)$ is the cocomponent segment of $\deleted{\unsure{A}}S_i$ which is contained in $A$, then the condition that $S_i,S_j$ are \deleted{\unsure{dis}}\unsure{not|indeed} mixed is expressed by the condition that $A_i\cap A_j=\phi$, hence $A_i,A_j$ are saturated segments $\subset A$ and $\ne A$, absurd.
\end{minipage}
\normalmarginpar
% archive sheet number: 48
